% !TeX root = ../../tesi.tex

\section{Motivazioni e obiettivi personali}
\label{sec:obiettivi-personali}

Da marzo 2026 lavoro in proprio come sviluppatore di negozi Shopify, occupandomi di migrazioni da
altre piattaforme, gestione dei negozi e automazioni. La mia esperienza riguardava però soltanto i
temi Liquid e non lo sviluppo di applicazioni. Tra le proposte di \emph{stage}, quella di Halue
riguardava progetti su Shopify, era la più vicina al mio lavoro e mi offriva ciò che mi mancava:
lo sviluppo di un’applicazione in un’azienda che lavora su questa piattaforma ogni giorno. Per
questo l’ho scelta.

L’idea dell’interfaccia generativa è di Halue. Quando l’azienda me l’ha proposta vi ho riconosciuto
l’esigenza di un mio cliente, che mi aveva parlato di un assistente capace di dare assistenza ai
visitatori del suo negozio e di indirizzarli verso l’acquisto. Il progetto aveva quindi per me un
valore che andava oltre lo \emph{stage}, e questo ha orientato sia la scelta del tema sia la
proposta dei componenti.

Da queste ragioni ho ricavato quattro obiettivi personali (\tabref{tbl:personali}), che valuto
insieme a quelli del piano nella \secref{sec:soddisfacimento}.

\begin{tabellatesi}{P{0.1}P{0.9}}
	{I miei obiettivi personali per lo \emph{stage}.}
	{tbl:personali}
	\toprule
	\rowcolor{TabellaIntestazione} \textbf{ID} & \textbf{Obiettivo} \\
	\midrule
	OP1 & Estendere le mie competenze su Shopify dai temi allo sviluppo di applicazioni \\
	OP2 & Acquisire esperienza nella progettazione di sistemi basati su agenti \\
	OP3 & Ricavare dal progetto conoscenze e soluzioni riutilizzabili nella mia attività
	professionale \\
	OP4 & Gestire in autonomia il lavoro da remoto, dove nessuno segue la mia giornata in modo
	continuo e spetta a me organizzarla, stabilire le priorità e prevedere le pause \\
	\bottomrule
\end{tabellatesi}

% DA CONFERMARE: il cliente resta anonimo finché non autorizza la citazione del nome.
